\fichatitulo{48 · AVALIAÇÃO DE CITAÇÕES}{INLG · 2024}{Suporte fino de citações}
{\small Zhang et al.. \textit{Suporte fino de citações}. INLG · 2024. \href{https://aclanthology.org/2024.inlg-main.35/}{Fonte consultada}.\par}

\campo{Problema e objetivo}
Uma citação sustenta integralmente a afirmação que acompanha?

\campo{Principal contribuição · síntese interpretativa}
Propõe avaliar suporte em três níveis: completo, parcial e ausente, comparando métricas de fidelidade.

\campo{Método / natureza da evidência}
Combina correlação, classificação e avaliação de recuperação contra julgamentos humanos.

\campo{Resultado ou argumento principal}
Nenhuma métrica domina consistentemente todos os testes; suporte parcial revela limites da classificação binária.

\campo{Excertos-chave · seleção literal}
\begin{itemize}
\excerto{1}{full, partial, and no support}{Resumo.}
\end{itemize}

\campo{Limitações e análise crítica}
O estudo avalia citações em geração textual e não garante que o esquema de três classes seja suficiente para afirmações legislativas complexas. Ainda assim, impede tratar presença de link como verificabilidade.

\campo{Aproveitamento na dissertação · proposta}
Metodologia: registrar suporte por afirmação e permitir a categoria parcial na rubrica.

\registro{Fonte consultada nesta preparação: PDF publicado; consulta focal do resumo, introdução, resultados e conclusão. Chave: \texttt{zhangEtAl2024fineGrainedCitation}.}
